% source: David Mumford, "Abelian Varieties", TIFR Studies in Mathematics 5, Oxford UP, 1970
% pdf_page: 0054
% doc_page: 43 (Chapter II, Section 4, "Definition of abelian varieties")
% retrieved: 2026-07-11
% transcribed_by: claude-fable-5 (vision, rendered at 200dpi via PyMuPDF; scanned book, no text layer)

% [running head: ALGEBRAIC THEORY VIA VARIETIES, page 43]

% [end of the proof that n_X is surjective, from the previous page]
It follows by induction on $n$ that for any $n > 0$ (and hence also for
$n \le 0$), $(dn_X)_0$ is multiplication by $n$. Thus, if $p \nmid n$, $dn_X$
is an isomorphism. If $n_X$ were not surjective, by the dimension theorem,
$\dim_0 n_X^{-1}(0) > 0$, and we can therefore find a non-zero $t \in T$
tangent to $n_X^{-1}(0)$ at $0$. But then, we would have $dn_X(t) = 0$
(since $n_X^{-1}(0)$ is mapped into the single point $0$ by $n_X$), which is
a contradiction.

The following lemma, besides having other important applications, gives a
second proof of the commutativity of $X$.

\paragraph{Rigidity Lemma. (Form I.)}
\textit{Let $X$ be a complete variety, $Y$ and $Z$ any varieties, and
$f : X \times Y \to Z$ a morphism such that for some $y_0 \in Y$,
$f(X \times \{y_0\})$ is a single point $z_0$ of $Z$. Then there is a
morphism $g : Y \to Z$ such that if $p_2 : X \times Y \to Y$ is the
projection, $f = g \circ p_2$.}

\paragraph{Proof.}
Choose any point $x_0 \in X$, and define $g : Y \to Z$ by
$g(y) = f(x_0, y)$. Since $X \times Y$ is a variety, to show that
$f = g \circ p_2$, it is sufficient to show that these morphisms coincide on
some open subset of $X \times Y$. Let $U$ be an affine open neighbourhood of
$z_0$ in $Z$, $F = Z - U$, and $G = p_2(f^{-1}(F))$; then $G$ is closed in
$Y$ since $X$ is complete and hence $p_2$ is a closed map. Further
$y_0 \notin G$ since $f(X \times \{y_0\}) = \{z_0\}$. Therefore
$Y - G = V$ is a non-empty open subset of $Y$. For each $y \in V$, the
complete variety $X \times \{y\}$ gets mapped by $f$ into the affine variety
$U$, and hence to a single point of $U$. But this means that for any
$x \in X$, $y \in V$, $f(x, y) = f(x_0, y) = g \circ p_2(x, y)$, and this
proves our assertion.

\paragraph{Corollary 1.}
\textit{If $X$ and $Y$ are abelian varieties and $f : X \to Y$ is any
morphism, $f(x) = h(x) + a$ where $h$ is a homomorphism of $X$ into $Y$ and
$a \in Y$.}

\paragraph{Proof.}
Replacing $f$ by $f - f(0)$, we may assume $f(0) = 0$ and we have to show
under this assumption that $f$ is a homomorphism. Consider the morphism
$\phi : X \times X \to Y$ defined by
$\phi(x, y) = f(x + y) - f(y) - f(x)$. Then
$\phi(X \times \{0\}) = \phi(\{0\} \times X) = 0$, so that it follows by the
above lemma that $\phi \equiv 0$ on $X \times X$, or what is the same, that
$f$ is a homomorphism.
